% Supplemental material (separate PDF). Compile after main.tex so xr can read main.aux.
\documentclass[letterpaper,journal]{IEEEtran}
%\usepackage{ctex}
\usepackage{amsmath,amsfonts}
\usepackage{amsthm}
\usepackage{algorithmic}
\usepackage{algorithm}
\usepackage{array}
\usepackage[caption=false,font=normalsize,labelfont=sf,textfont=sf]{subfig}
\usepackage{textcomp}
\usepackage{stfloats}
\usepackage{url}
\usepackage{verbatim}
\usepackage{graphicx}
\usepackage{epsfig}
% \usepackage{subfigure}
\usepackage{epstopdf}
\usepackage{cite}
\usepackage{booktabs}
\usepackage{multirow} 
\usepackage{amssymb} 
\newtheorem{definition}{Definition}  
\newtheorem{theorem}{Theorem}
\newtheorem{proposition}{Proposition}
\newtheorem*{proposition*}{Proposition}

\newtheorem{assumption}{Assumption}
\newtheorem{lemma}{Lemma}
\newtheorem{corollary}{Corollary}
\newtheorem{remark}{Remark}
% End-of-proof marker. \qed ends the paragraph and sets the box flush right on
% its own line, so the sentence that follows a proof starts as a new paragraph.
% (amsthm defines \qedhere but not a top-level \qed.)
\providecommand{\qed}{\hfill$\square$\par}
%\usepackage{appendix}
\usepackage{placeins}
\graphicspath{{./figures/}{./images/}}
\hyphenation{op-tical net-works semi-conduc-tor IEEE-Xplore}

\makeatletter
\newcommand{\rmnum}[1]{\romannumeral #1}
\newcommand{\Rmnum}[1]{\expandafter\@slowromancap\romannumeral #1@}
\makeatother

\renewcommand{\topfraction}{0.99}
\renewcommand{\bottomfraction}{0.99}
\renewcommand{\textfraction}{0.01}
\renewcommand{\floatpagefraction}{0.99}
\renewcommand{\dbltopfraction}{0.9}
\renewcommand{\dblfloatpagefraction}{0.9}
\setlength{\dbltextfloatsep}{8pt plus 2pt}
\setcounter{topnumber}{5}
\setcounter{bottomnumber}{5}
\setcounter{totalnumber}{10}


\usepackage{xr}
\externaldocument{main}

\begin{document}

\title{Supplemental Material for\\
RAVEN-MCS: Target-Mass-Calibrated Windowed-Asynchronous Federated Data Completion for Mobile Crowdsensing under Two-Stage Non-Random Selection}

\author{**** and Anfeng Liu}

\maketitle

\noindent\textit{This file is supplemental material for the companion main paper
and must be uploaded as a separate PDF. It is not a standalone submission.}

% Continue figure/table/equation numbers from the companion main paper.
% Main now has 4 figures, 4 tables, 1 algorithm (window-close Design rule) and
% the support-gate / switch-point equations (42)-(43).
% Numbered statements in main: Prop. 1--3, Thm. 1--2, Assumption 1.
\setcounter{figure}{4}
\setcounter{table}{4}
\setcounter{algorithm}{1}
\setcounter{equation}{43}
\setcounter{proposition}{4}
\setcounter{theorem}{2}
\setcounter{corollary}{0}
\setcounter{assumption}{1}

\input{appendix}

% Trained-rule probe deferred from the results section of the main paper.  It is
% placed after \input{appendix} so that the appendix keeps its own section
% numbering, and it adds no numbered statement, so the counters set above stay
% valid for the appendix.
\section{Exploratory Trained-Rule Probe}
\label{app:trained_rule_probe}

\subsection{The $\kappa=1000$ projection-with-debt rule}
\label{app:kappa_probe}

The projection-with-debt rule is a new rule of this study, not the frozen RAVEN
configuration, and it is reported as exploratory: it sets the debt weight to
$\kappa=1000$ instead of the frozen value, and that weight was chosen on the
tuning seeds.  Against the projection alone it adds $+1.347$ pp and $+0.580$ pp
(Holm $\approx10^{-5}$ on both), on SensorScope and U-Air respectively; at the
frozen debt strength the same contrast is $+0.0073$ pp ($[0.0064,0.0083]$) on
SensorScope, inside the equivalence band of
Section~\ref{subsec:equivalence_band} in the main paper.  Its implementation
gate is not met.  The trained rules are reported as approximate against an exact
projection oracle, measured by $\varepsilon_{\mathrm{impl}}$, the distance
between a trained rule's server weights and the exact minimizer of its
projection problem, in percentage points of group mass rather than as a
relative tolerance.  Over all $240$ rule runs
$\varepsilon_{\mathrm{impl}}\le2.42\times10^{-5}$ percentage points, at most
$0.11\%$ of any of the four Inst readings and $2.16\%$, $7.62\%$, $2.73\%$ and
$4.89\%$ of the four debt readings.  The two U-Air debt readings are the
exception, where the gap exceeds the reading it is compared against.  Elsewhere
it is orders of magnitude below the effect, so no sign or reading changes.
These deltas belong to the new rule, and none of them is a measurement of
RAVEN's performance.

% The appendix cites a small number of works. Carrying its own bibliography
% keeps this file compilable on its own, without main.aux.
\bibliographystyle{IEEEtran}
\bibliography{reference}

\end{document}
